\documentclass[12pt,a4paper]{article}
\usepackage[utf8]{inputenc}
\usepackage[T1]{fontenc}
\usepackage[french]{babel}
\usepackage{geometry}
\geometry{a4paper, margin=2.5cm, top=2.5cm, bottom=2.5cm}
\usepackage{setspace}
\setstretch{1.25}
\usepackage{hyperref}
\usepackage{listings}
\usepackage{xcolor}
\usepackage{booktabs}
\usepackage{array}
\usepackage{fancyhdr}
\usepackage{mdframed}

\pagestyle{fancy}
\fancyhf{}
\rhead{\textbf{Compte-Rendu Technique : VBA, Référencement \& VLOOKUP}}
\lhead{Ingénierie de Données Excel}
\rfoot{Page \thepage}

\definecolor{codegreen}{rgb}{0,0.6,0}
\definecolor{codegray}{rgb}{0.5,0.5,0.5}
\definecolor{codepurple}{rgb}{0.58,0,0.82}
\definecolor{backcolour}{rgb}{0.96,0.96,0.94}
\definecolor{primaryblue}{rgb}{0.1,0.3,0.6}

\lstdefinestyle{mystyle}{
    backgroundcolor=\color{backcolour},   
    commentstyle=\color{codegreen},
    keywordstyle=\color{primaryblue}\bfseries,
    numberstyle=\tiny\color{codegray},
    stringstyle=\color{codepurple},
    basicstyle=\ttfamily\footnotesize,
    breakatwhitespace=false,         
    breaklines=true,                 
    captionpos=b,                    
    keepspaces=true,                 
    numbers=left,                    
    numbersep=7pt,                  
    showspaces=false,                
    showstringspaces=false,
    showtabs=false,                  
    tabsize=2
}
\lstset{style=mystyle}

\title{\Huge \textbf{Rapport de Réunion Technique}\\[0.8em]
\Large Automatisations VBA, Syntaxe de Référencement Inter-Feuilles (\$, !) et Pratiques de la Fonction VLOOKUP (RECHERCHEV)}
\author{\textbf{Département des Systèmes d'Information \& Ingénierie des Données}}
\date{\today}

\begin{document}
\maketitle

\begin{mdframed}[backgroundcolor=backcolour,linecolor=primaryblue,linewidth=1.5pt]
\textbf{Fiche Synthétique de la Réunion}
\begin{itemize}
    \item \textbf{Objet :} Standardisation des pratiques d'automatisation Excel, syntaxe de liaison inter-feuilles et requêtes VLOOKUP.
    \item \textbf{Date \& Heure :} \today{} -- 10h00 à 12h00
    \item \textbf{Lieu :} Salle de Conférence DSI / Visioconférence
    \item \textbf{Rédacteur :} Équipe d'Ingénierie des Données
    \item \textbf{Statut du Document :} Compte-Rendu Officiel Validé
\end{itemize}
\end{mdframed}

\vspace{1em}
\tableofcontents
\newpage

\section{Introduction et Ordre du Jour}
Cette réunion technique s'est tenue dans le cadre du projet d'optimisation et de sécurisation des classeurs de travail partagés au sein de l'organisation. L'objectif principal était de clarifier les règles d'automatisation par macros VBA, de poser des standards stricts sur les références inter-feuilles utilisant les opérateurs \texttt{!} et \texttt{\$}, et d'établir un guide de bonnes pratiques sur l'utilisation de la fonction de recherche verticale \texttt{VLOOKUP} (\texttt{RECHERCHEV}).

\section{Axe 1 : Le Moteur VBA et l'Automatisation des Traitements}

\subsection{Principes d'Intégration du Moteur VBA}
Le Visual Basic for Applications (VBA) constitue le moteur d’automatisation historique d’Excel. Les points d'architecture suivants ont été rappelés :
\begin{itemize}
    \item \textbf{Environnement d'Édition (VBE) :} Accessible via le raccourci \texttt{Alt + F11} ou par l'onglet \textit{Développeur}.
    \item \textbf{Modules Standards vs Modules Événementiels :} 
    \begin{itemize}
        \item Les modules standards (\texttt{Module1}) hébergent les procédures globales exécutables à la demande via des boutons.
        \item Le module \texttt{ThisWorkbook} héberge le code réactif s'exécutant automatiquement sur les événements système tels que \texttt{Workbook\_SheetChange} ou \texttt{Workbook\_Open}.
    \end{itemize}
\end{itemize}

\subsection{Exemple de Macro Événementielle Inter-Feuilles}
Le script ci-dessous a été présenté comme standard pour intercepter en temps réel la modification d'une feuille maître et propager les valeurs vers une feuille copie :

\begin{lstlisting}[caption=Macro événementielle de synchronisation inter-feuilles]
Private Sub Workbook_SheetChange(ByVal Sh As Object, ByVal Target As Range)
    Const MASTER_NAME As String = "FeuilleMaster"
    Dim wsCopy As Worksheet
    
    ' Ne traiter que les modifications sur la feuille master
    If Sh.Name <> MASTER_NAME Then Exit Sub
    
    On Error GoTo CleanUp
    Application.EnableEvents = False
    
    For Each wsCopy In ThisWorkbook.Worksheets
        If wsCopy.Name <> MASTER_NAME Then
            wsCopy.Range(Target.Address).Value = Target.Value
        End If
    Next wsCopy

CleanUp:
    Application.EnableEvents = True
End Sub
\end{lstlisting}

\subsection{Gouvernance et Sécurité des Fichiers avec Macros}
Pour garantir la sécurité et prévenir les blocages par les pare-feux d'entreprise :
\begin{enumerate}
    \item Les classeurs contenant des macros doivent obligatoirement être enregistrés au format \textbf{\texttt{.xlsm}} (ou \textbf{\texttt{.xlsb}}).
    \item Les projets VBA déployés en production doivent être signés numériquement à l'aide d'un certificat d'entreprise reconnu par le Centre de confiance.
\end{enumerate}

\section{Axe 2 : Syntaxe du Référencement Inter-Feuilles (Symboles ! et \$)}

\subsection{L'Opérateur de Séparation de Feuille : Le Point d'Exclamation (!)}
Pour faire référence à une cellule située dans une autre feuille du même classeur, Excel exige l'utilisation du point d'exclamation \texttt{!} comme séparateur entre le nom de la feuille source et la plage de cellules.

\begin{itemize}
    \item \textbf{Syntaxe standard (Nom de feuille simple) :} \\
    \texttt{=FeuilleMaster!A1} \\
    Cette formule récupère la valeur de la cellule A1 dans la feuille nommée \textit{FeuilleMaster}.
    \item \textbf{Syntaxe avec espaces ou caractères spéciaux :} \\
    Si le nom de la feuille comporte des espaces, il doit être impérativement encadré par des guillemets simples \texttt{'} : \\
    \texttt{='Donnees Ventes'!A1}
\end{itemize}

\subsection{L'Opérateur de Verrouillage de Référence : Le Signe Dollar (\$)}
Le signe dollar \texttt{\$} permet d'ancrer les coordonnées d'une cellule afin d'éviter qu'elles ne se décalent lors de l'étirement (recopie vers le bas ou vers la droite) de la formule.

La réunion a formalisé les quatre types de référencement :

\begin{table}[h]
\centering
\begin{tabular}{@{}lll@{}}
\toprule
\textbf{Type de Référence} & \textbf{Exemple Syntaxe} & \textbf{Comportement lors de la Recopie} \\ \midrule
Relative & \texttt{FeuilleMaster!A1} & Ajuste la ligne et la colonne \\
Absolue & \texttt{FeuilleMaster!\$A\$1} & Ligne et colonne strictement bloquées \\
Mixte (Ligne fixée) & \texttt{FeuilleMaster!A\$1} & La colonne varie, la ligne reste fixée à 1 \\
Mixte (Colonne fixée) & \texttt{FeuilleMaster!\$A1} & La colonne reste fixée à A, la ligne varie \\ \bottomrule
\end{tabular}
\caption{Matrice des types de référencement inter-feuilles avec \$}
\end{table}

\subsection{Cas Pratique : Étirement d'une Formule Inter-Feuilles}
Lorsqu'une formule de consolidation inter-feuilles est étirée sur un grand tableau :
\begin{itemize}
    \item Si on écrit \texttt{='Master'!\$B\$2}, la recopie sur 1000 lignes renverra toujours la valeur de B2 (Erreur fréquente).
    \item Si on écrit \texttt{='Master'!B2}, la recopie s'adaptera proprement en B3, B4, B5, etc.
    \item Si on pointe sur une table de référence fixe, l'utilisation de \texttt{='Master'!\$A\$1:\$D\$100} est obligatoire pour empêcher la plage de glisser vers le bas.
\end{itemize}

\section{Axe 3 : Maîtrise de la Fonction VLOOKUP (RECHERCHEV)}

\subsection{Structure et Syntaxe Officielle}
La fonction \texttt{VLOOKUP} (ou \texttt{RECHERCHEV} en version française) permet de rechercher une valeur dans la première colonne d'une table et de renvoyer une valeur située sur la même ligne dans une autre colonne.

\textbf{Syntaxe :} \\
\texttt{=VLOOKUP(valeur\_cherchee; table\_recherche; index\_colonne; [correspondance\_exacte])}

\begin{enumerate}
    \item \texttt{valeur\_cherchee} : La clé de recherche (ex: un ID client ou un code produit).
    \item \texttt{table\_recherche} : La plage de données dans laquelle effectuer la recherche.
    \item \texttt{index\_colonne} : Le numéro de la colonne à partir de laquelle la valeur doit être extraite (1 pour la première colonne).
    \item \texttt{correspondance\_exacte} : Mettre \texttt{FALSE} (ou \texttt{0}) pour exiger une correspondance exacte.
\end{enumerate}

\subsection{Association de VLOOKUP avec le Référencement Inter-Feuilles}
Dans les architectures d'entreprise, la table de référence réside presque toujours dans une feuille séparée (ex: feuille \textit{Catalogue}).

\textbf{Formule de Référence Validée :}
\begin{quote}
\texttt{=VLOOKUP(A2; 'Catalogue'!\$A\$2:\$D\$1000; 3; FALSE)}
\end{quote}

\textbf{Analyse des éléments de cette formule :}
\begin{itemize}
    \item \texttt{A2} : Valeur recherchée dans la feuille courante (référence relative).
    \item \texttt{'Catalogue'!\$A\$2:\$D\$1000} : Référencement inter-feuilles avec le séparateur \texttt{!} et ancrage absolu \texttt{\$} pour éviter tout décalage de la matrice lors de la recopie vers le bas.
    \item \texttt{3} : Renvoie le tarif situé dans la 3ème colonne du tableau catalogue.
    \item \texttt{FALSE} : Garantit qu'aucun résultat erroné ne sera renvoyé si le code produit n'existe pas.
\end{itemize}

\subsection{Limites Techniques de VLOOKUP et Alternatives Recommandées}
Les participants ont identifié deux contraintes majeures de \texttt{VLOOKUP} :
\begin{enumerate}
    \item \textbf{Sensibilité à l'insertion de colonnes :} Si une colonne est insérée au milieu du tableau source, l'index numérique fixe (\texttt{3}) renvoie une mauvaise colonne.
    \item \textbf{Impossible de chercher à gauche :} La clé de recherche doit obligatoirement se trouver dans la première colonne de la plage.
\end{enumerate}

\textbf{Solutions modernes préconisées par la DSI :}
\begin{itemize}
    \item La combinaison \texttt{INDEX} + \texttt{MATCH} (\texttt{INDEX} + \texttt{EQUIV}).
    \item La fonction \texttt{XLOOKUP} (\texttt{RECHERCHEX}) disponible sous Microsoft 365, plus rapide et insensible à l'inséraction de colonnes.
\end{itemize}

\section{Bilan des Décisions et Recommandations}

\begin{table}[h]
\centering
\begin{tabular}{@{}llp{6cm}@{}}
\toprule
\textbf{N°} & \textbf{Sujet} & \textbf{Décision / Recommandation} \\ \midrule
1 & Format des Fichiers & Généraliser l'usage du format \texttt{.xlsm} pour tous les fichiers avec macros VBA. \\
2 & Ancrage \$ & Exiger le verrouillage absolu \texttt{\$A\$1:\$Z\$100} pour toutes les tables sources inter-feuilles. \\
3 & Recherche & Imposer le paramètre \texttt{FALSE} (0) dans toutes les fonctions \texttt{VLOOKUP}. \\
4 & Migration M365 & Encourager le passage progressif de \texttt{VLOOKUP} vers \texttt{XLOOKUP}. \\ \bottomrule
\end{tabular}
\caption{Synthèse des décisions de la réunion}
\end{table}

\section{Conclusion}
Les normes validées lors de cette réunion permettront d'éliminer les erreurs récurrentes de décalage de formules et d'assurer une meilleure stabilité des classeurs partagés au sein de l'organisation.

\end{document}
